\section{Related Work}
\label{sec:related}

\tool is inspired by work on library learning,
specifically the \dc line of work,
as well as equality saturation-based program synthesis and decompilation.

\mypara{\dc}
%
\dc~\cite{dreamcoder} is a program synthesizer
  that learns a library of abstractions
  from solutions to a set of related synthesis tasks.
The library is intended to be used for solving other
  similar synthesis tasks.
% The DSLs for the different domains in \dc
%   are carefully designed to help scale the search.
\dc
  uses version spaces~\cite{vsa1, vsa2}
  to compactly store a large
  set of programs
  and leverages ideas from \egraphs
  (such as e-matching)
  but only for exploring the space of refactorings
  of the original program using the candidate libraries,
  not for making library learning robust to syntactic variation.
%
Our evaluation shows that \tool can find more optimal abstractions faster than \dc.

\dc has sparked several direction of follow-up work
that attempt to improve the efficiency of its library learning procedure
and the quality of the learned abstractions.
%
One of them is by \citet{laps},
which uses natural language annotations and a neural network
to guide library learning.
%
Another one is \stitch~\cite{stitch},
which was developed \emph{concurrently} with this work
and is the most closely related to \tool;
we discuss \stitch in some detail below.

\mypara{\stitch}
%
The core difference between the two approaches
is that \stitch focuses on improving \emph{efficiency} of purely syntactic library learning,
whereas \tool attempts to improve its \emph{expressiveness} by adding equational theories.
%
While \tool separates library learning into two phases---%
candidate generation via anti-unification and candidate selection via e-graph extraction---%
the \stitch algorithm \emph{interleaves} the generation and selection phases
in a branch-and-bound top-down search.
%
Starting from the ``top'' pattern $X$,
\stitch gradually refines it until further refinement does not pay off.
%
To quickly prune suboptimal candidate patterns,
\stitch computes an upper bound on their compression
by summing up the compression at each match of this pattern in the corpus
(this bound is imprecise because it does not take into account that matches might overlap).
%
For candidates that are not pruned this way,
\stitch computes their true compression by searching for the optimal subset of matches to rewrite,
a so-called ``rewrite strategy''.
%
\tool's extraction algorithm can be seen as a generalization of \stitch's rewrite strategy:
while the former searches over both subsets of patterns and how to apply them to the corpus at the same time,
the latter considers a single pattern at a time and only searches for the best way to apply it.
%
Since the search space in the former case is much larger,
\tool uses a beam search approximation,
while in \stitch the rewrite strategy is precise.
%
To sum up, the main pros and cons of the two approaches are:
\begin{itemize}
\item \tool can learn libraries modulo equational theories, while \stitch cannot;
\item \stitch provides optimality guarantees for learning a single best abstraction at a time, while \tool can learn multiple abstractions at once, but sacrifices theoretical optimality.
\end{itemize}


\mypara{Other Library Learning Techniques}
%
\tname{Knorf}~\cite{knorf} is a library learning tool for logic programs,
which, like \tool, proceeds in two phases.
%
Their candidate generation phase is similar to the upper bound computation in \stitch,
while their selection phase uses an off-the-shelf constraint solver.
%
It would be interesting to explore whether their constraint-based technique can be generalized beyond logic programs.

Other work develops limited forms of library learning,
where only certain kinds of sub-terms can be abstracted.
%
For example, ShapeMod~\cite{shapemod} learns macros
  for 3D shapes represented
  in a DSL called ShapeAssembly,
  and only supports abstracting over numeric parameters,
  like dimensions of shapes.
%
Our own prior work~\cite{cogsci-smiley} extracts common structure from graphical programs,
  but only supports abstracting over primitive shapes
  and applying the abstraction at the top level of the program.
%
Such restrictions make the library learning problem more computationally tractable,
  but limit the expressiveness of the learned abstractions.

There are several neural program synthesis tools%
  ~\cite{polozov, srini, learner, gredilla}
  that learn \emph{programming idioms}
  using statistical techniques.
  % learn programming idioms that make synthesis more efficient.
Some of these tools have used
  ``explore-compress'' algorithms~\cite{learner} to
  iteratively enumerate a set of
  programs from a grammar and find a solution
  that exposes abstractions that make the set of
  programs maximally compressible.
This is similar to
  common subexpression elimination which \tool uses
  for guiding extraction.
%\cite{gredilla} have used a similar approach for
%  helping robots learn to complete various tasks by
%  selecting from a library of abstractions.
%A recent work from us combines this
%  explore-compress algorithm with \egraphs
%  for co-optimizing designs and fabrication instructions
%  for carpentry~\cite{cc2}
%  although it does not focus on synthesizing
%  new abstractions.

% \tool's library learning is based on \egraph
%   anti-unification combined with DSRs, which differs
%   different from prior approaches; our
%   results show that it can find better compressions
%   faster than the state-of-the-art~\cite{dreamcoder}.


\mypara{Loop rerolling}
%
Loop rerolling is related to library learning in that it also aims to discover hidden structure in a program,
except that this structure is in the form of loops.
%
A variety of domains have used loop rerolling
  to infer abstractions from flat input programs.
In hardware, loop rerolling is
  used to optimize programs
  for code size~\cite{hardware1, hardware2, hardware3}.
In many of these tools,
  the compiler first unrolls a loop,
  applies optimizations, then rerolls it --- the compiler
  therefore has structural information about the loop
  that can be used for rerolling~\cite{hardware1}.
The graphics domain has used
  loop-rerolling to discover
  latent structure from low-level representations.
CSGNet~\cite{csgnet} and Shape2Prog~\cite{shape2prog} used
  neural program generators to discover for loops from
  pixel- and voxel-based input representations.
\cite{handdrawn} used program synthesis
  and machine learning
  to infer loops from hand-drawn images.
Szalinski~\cite{szalinski} used
  equality saturation to
  automatically learn loops in the form
  of maps and folds
  from flat 3D CAD programs that are
  synthesized by mesh decompilation tools~\cite{reincarnate}.
WebRobot~\cite{webrobot} has used speculative
  rewriting for inferring loops from traces
  of web interactions.
% WebRobot is similar to Szalinski in that they both
%   use the notion of speculative rewrites.
Similar to \tool (and unlike Szalinski),
  WebRobot finds abstractions
  over \textit{multiple} input traces.
% Compared to these tools however,
%   \tool's approach is more general --- it is not
%   restricted to finding only loops or limited
%   types of macros and does not rely on
%   any information
%   about how the original
%   ``flat'' programs were generated.

\mypara{Applications of Anti-Unification}
%
Anti-unification is a well-established technique for discovering common structure in programs.
%
It is the core idea behind bottom-up Inductive Logic Programming~\cite{ilp},
and has also been used for software clone detection~\cite{Bulychev10},
programming by example~\cite{raza2014programming},
and learning repetitive code edits~\cite{lase,refaser}.
%
It is possible that these applications could also benefit from \tool's notion of anti-unification over e-graphs
to make them more robust to semantics-preserving transformations.

\mypara{Synthesis and Optimization using \Egraphs}
While traditionally \egraphs have been used
  in SMT solvers for
  facilatiting communication between
  different theories,
  several tools have demonstrated
  their use for optimization
  and synthesis.
\citet{tate2009equality} first used \egraphs for
  equality saturation: a rewrite-driven technique
  for optimizing Java programs with loops.
Since then, several tools have used equality saturation
  for finding programs equivalent to,
  but better than, some
  input program~\cite{willsey2021egg, herbie, tensat, szalinski, diospyros, spores,
  cc1}.
\tool uses an anti-unification
  algorithm on \egraphs
  (together with domain specific
  rewrites), which prior work has not shown.
Additionally, prior work has either used
  greedy or ILP-based extraction strategies,
  whereas \tool uses a new targeted common
  subexpression elimination approach which
  we believe can be used in many other applications of
  equality saturation, especially given its
  amenability to approximation via beam search.


